\documentclass[journal]{IEEEtran}
\usepackage{newtxtext,newtxmath}
\usepackage{booktabs}
\usepackage{array}
\usepackage{amsmath}
\usepackage{graphicx}
\usepackage[hidelinks]{hyperref}

\newcommand{\modeltag}{\textsuperscript{[model]}}
\newcommand{\vcstag}{\textsuperscript{[VCS]}}
\newcommand{\dctag}{\textsuperscript{[DC/PT]}}
\newcolumntype{L}[1]{>{\raggedright\arraybackslash}p{#1}}

\markboth{IEEE Transactions on Circuits and Systems---II: Express Briefs}{Authors: Finite-Capacity Exact Product Capture}

\title{Exact Product Capture under a Finite Single-Port Parent Store\\
for Event-Driven Spiking Optical Flow}

\author{Authors to be inserted%
\thanks{Manuscript received Month Day, 2026; revised Month Day, 2026.}%
\thanks{The authors are with Huazhong University of Science and Technology,
Wuhan 430074, China. TCAS-II is single-blind: replace this block with legal
names, IEEE membership, e-mail, ORCID, and funding before submission.}}

\begin{document}
\maketitle

\begin{abstract}
Low firing activity does not automatically accelerate event-driven optical
flow when parent-state traffic is bound by finite storage and a single 1RW
port.  This brief reports one exact product-capture island, C1, for a frozen
binary-event optical-flow checkpoint.  A row may reuse only a still-resident
exact-subset parent, issues the XOR residual, and reconstructs the signed
product row by addition.  Dead writes that no later row can observe are
elided.  Completion of a row and its final psum write are atomic under a
deadline-aware 1RW arbiter.  On 51.84 million source rows from ten
\texttt{zurich\_city\_09\_a} samples, the same-ledger cycle model reduces
component time by 40.99\% (1.6945$\times$) versus strongest-zero skipping;
same-coordinate bit skipping on that ledger is only 1.003$\times$.  The mapped
28-nm nine-SRAM island occupies 166{,}514~$\mu$m$^2$, meets 3-ns setup/hold,
dissipates 29.08~mW / 22.07~nJ in a 64-row/253-cycle directed window, and
passes 16{,}549 mapped-to-mapped Formality compare points.  A real-mask VCS
tile calibrates event counts; the 1.6945$\times$ ratio remains a cycle-model
result.  All physical numbers are prelayout and pre-macro-extracted.  We make
no whole-network speedup claim.
\end{abstract}

\begin{IEEEkeywords}
Digital accelerator, event-based vision, optical flow, product sparsity,
single-port SRAM, spiking neural networks.
\end{IEEEkeywords}

\IEEEPARstart{E}{vent} cameras make inactivity visible, but sparse arithmetic
is only one term in execution time.  In a binary-event optical-flow
deployment, nonzero sources still trigger parent-state reads, residual
issues, and psum writes.  A design that reports only operation sparsity can
therefore be slower than an equal-resource zero-skipping baseline once the
parent working set is finite and the store is single-ported.

Prosperity~\cite{prosperity} and Phi~\cite{phi} show that useful SNN speedups
come from executable reuse or patterns, not firing rate alone.  Digital
neuromorphic processors such as ODIN~\cite{frenkel2019odin} and
event-stream accelerators such as SpinalFlow and SNE~\cite{spinalflow,sne}
establish spike-driven datapaths.  TCAS-II SRAM compute-in-memory macros
exploit sparsity inside the array~\cite{qiao2022dscim,taylor2022pim}.  Our
object is different: repeated \emph{products} of a convolution-like layer,
subject to a 240-KiB-class parent lifetime and a physical 1RW contract.

This brief makes one architectural contribution.  C1 captures exact repeated
products while a legal parent remains resident, reconstructs the child from
the residual, and never exposes a partial row.  The frozen Motion checkpoint
is an evaluation workload, not a parallel algorithmic claim.  A companion
typed weight-broadcast fabric for fully connected layers exists in the same
project; its hold and power are not closed, so it is not claimed here.

\section{Workload and Exactness Contract}
The frozen Motion C12 ep34 checkpoint obtains 1.199514 average endpoint error
(AEE) on an 825-frame local DSEC validation set and has a 5.6709\% global
firing rate.  Of 105 configured ATLIF wrappers, 12 \texttt{sn2\_q} wrappers
are runtime-bypassed, so capture invokes 93 (48 $T{=}2$, 45 $T{=}10$).  A
graph audit finds that 12 invoked \texttt{attn\_sn} returns have no consumer
under fixed normal inference, leaving 81 graph-live services.  All 93
captured ATLIF outputs are binary.  Signed source codes are a downstream
execution-protocol property (polarity and correction traffic), not evidence
that the captured tensors are analog.  The frozen software path uses
Motion-XOR scoring with $\alpha{=}0.125$ and K as its value carrier;
hardware quantization is disabled on that path.

Each result is assigned one evidence class.  A cycle-model result replays the
same trace and charged resources in software.  A VCS result proves directed
RTL function and protocol properties.  DC/PT/Formality results anchor area,
timing, and equivalence.  Only a result that closes all three classes on one
workload may be called an RTL speedup; this brief reports the three classes
side by side.

Bottleneck convolution exposes repeated products only while parent and
destination state remain resident.  That invariant defines C1.  Attention and
decoder traces remain part of the frozen identity but are not promoted to
performance contributions.

\begin{figure}[t]
\centering
\includegraphics[width=0.99\columnwidth]{figures/typed_source_islands.pdf}
\caption{Typed-source context of the frozen optical-flow network.  This brief
reports only C1: exact product capture in a finite 1RW parent store.  The
other branches are workload context and are not claimed.}
\label{fig:overview}
\end{figure}

\section{C1 Microarchitecture}
\subsection{Directory and Legal Parents}
C1 first builds a ping-pong directory for each 64-row, 16-source mask task.
A row may select only an earlier exact-subset row as its parent.  Stable
population-count ordering makes that choice unique: among legal parents the
engine prefers the earliest still-live candidate so that later rows keep a
deterministic residual.  A live-parent bitmap records which completed rows
remain observable.  When the bitmap shows that no subsequent row can read a
stored psum, that location is a dead write and is elided exactly rather than
being scheduled as a drain.

The 240-KiB-class working set is the capacity that must be respected, not a
DRAM-backed scratchpad.  Nine $128\times128$ 1RW SRAM macros hold parent
masks and psums.  The engine never assumes the ideal simultaneous 1R1W access
that produced a higher but unphysicalized reuse ceiling.  This is the object
difference from an unbounded product-sparsity opportunity: reuse must survive
capacity, port, parent lifetime, and completion constraints.

\subsection{Residual Issue and Exact Reconstruct}
A selected parent mask $\mathcal{M}_p\subset\mathcal{M}_r$ yields a disjoint
residual.  C1 issues only that residual and reconstructs the 96-lane signed
product row by adding the stored parent:
\begin{equation}
\mathbf{P}_r=\mathbf{P}_p+
\sum_{s\in\mathcal{M}_r\setminus\mathcal{M}_p}\mathbf{p}_s.
\label{eq:c1exact}
\end{equation}
Exactness follows from the disjoint union, not from value prediction.  If no
resident legal parent exists, C1 selects the empty parent and issues the
original source row.  Sign, destination, and accumulator identity are never
coalesced across rows.  Thus sparsity changes delivered work without changing
the operator result.

\subsection{Single-Port Arbiter and Atomic Commit}
The execution side combines one earliest-parent lookahead, a two-entry
reserved response queue, and forwarding under a deadline-aware 1RW arbiter.
A parent read, a residual product, and a psum write compete for the same
macro port.  The arbiter grants the access that would otherwise miss its
row-completion deadline; the reserved queue absorbs one already-accepted
response so that a later grant cannot orphan it.  Forwarding supplies a
just-written parent to the next residual without a second macro read when
the bitmap still marks that parent live.

Final psum write and row completion are a single atomic event.  Backpressure
may change service order, but it cannot expose a partially reconstructed row:
the typed-source terminal bit and the completion pulse retire together.  A
directed VCS tile on the first 64 real support masks of the cycle ledger is
used later only to calibrate these event counts, not to replace the
same-ledger cycle model.

Speculative parent hits are disallowed.  A candidate is legal only if its
mask is an exact subset and the live-parent bitmap is still set at grant
time, not at directory build time.  That second check is what makes
(1) exact under eviction: a parent that died between directory construction
and issue is treated as missing, and C1 falls back to the empty parent.
The fallback is slower than a stale hit and is the price of never
reconstructing from a psum that the store no longer owns.

\section{Evaluation}
All physical results use a commercial 28-nm flow at 3.0~ns.  Reported points
are prelayout, use ideal clocks and ZeroWireload, and have no extracted
parasitics.  C1 area includes the nine SRAM macros.  Logic-only comparisons
from other islands in the same project are not mixed into Table~\ref{tab:main}.

Table~\ref{tab:c1ladder} exposes the baseline ladder instead of attributing
all inactivity to product capture.  The ledger is 51.84 million source rows
from ten \texttt{zurich\_city\_09\_a} ep34 samples.  Per-sample ratios span
1.661--1.723$\times$, with a 1.6946$\times$ geometric mean.  Strongest-zero
skipping still walks parent-state traffic for every nonzero source;
same-coordinate bit skipping on this ledger removes almost none of that
traffic (1.003$\times$).  Exact parent reuse is what creates the material
reduction.  The remaining 12.25\% gap from C1 to the concurrent-access
ceiling is the tax of the physical single-port contract.  We do not promote
the ceiling as an attainable RTL result.

The cycle model charges the same resources the RTL would charge: one 1RW
grant per accepted macro access, residual issue, forward, deadline hold, and
atomic commit.  It is not a CPU-time measurement and is not an RTL speedup.
Only a future campaign that closes model, VCS, and DC/PT on one full ledger
may use that stronger label.

\begin{table}[t]
\caption{C1 same-ledger capture ladder (51.84~M source rows).}
\label{tab:c1ladder}
\centering
\scriptsize
\begin{tabular}{@{}lrrl@{}}
\toprule
Execution point & Cycles (M) & vs.\ zero & Boundary \\
\midrule
Strong zero & 648.741 & 1.000$\times$ & baseline \\
Same-coordinate bit & 646.619 & 1.003$\times$ & baseline \\
C1 finite 1RW & 382.849 & 1.6945$\times$ & cycle model \\
Concurrent-access & 341.058 & 1.902$\times$ & ceiling only \\
\bottomrule
\end{tabular}
\end{table}

The mapped nine-SRAM island occupies 166{,}514~$\mu$m$^2$.  Synthesis and
place-equivalent mapping use the foundry 28-nm standard-cell library and the
nine compiler SRAM abstracts; no Yosys or educational library is used as an
ASIC PPA proxy.  PrimeTime reports setup/hold slack of $+27.9$/$+1.8$~ps at
3~ns under ideal clocks.  Formality passes 16{,}549 mapped-to-mapped compare
points on that netlist pair.  In a separately bounded 64-row, 253-cycle
directed window the island dissipates 29.08~mW and 22.07~nJ.  That energy is
a mixed-corner, prelayout PrimeTime PX estimate without SPEF and without a
frame-level activity factor; it is not energy per DSEC frame and is not a
chip measurement.

A separately sealed ep34 tile uses the first 64 real support masks of the
cycle ledger.  VCS matches 196 accepted issues, 58 parent edges, 31
dead-write elisions, 54/33 macro reads/writes, four forwards, six deadline
holds, 14 stalls, and 64 commits/completions.  Lane values are deterministic
synthetic signed-12-bit data with zero prior psum.  This calibrates modeled
event counts and directed arithmetic; it does not itself produce the
1.6945$\times$ cycle ratio.  Protocol errors, X-pessimism on energy, and
failed sibling attempts in the same repository are not cited.

The 64-row task is the scheduling quantum that matches the ping-pong
directory, not a claim that the deployed convolution is only 64 rows deep.
Each task carries up to 16 typed sources and writes a 96-lane signed product.
The nine macros are addressed as a single 1RW parent store: one grant moves
either a mask, a psum, or a residual payload.  Capacity eviction follows the
live-parent bitmap rather than LRU of raw addresses, because a recently
written parent may already be dead while an older parent remains legally
observable.  That replacement rule is why the concurrent-access ceiling in
Table~\ref{tab:c1ladder} is not a drop-in RTL mode: it would require a second
port that the mapped island does not have.

The same-ledger model therefore reports two numbers that must not be
conflated.  1.6945$\times$ is C1 versus strongest-zero on the physical 1RW
contract.  1.902$\times$ is the unphysicalized ceiling.  The 12.25\% gap is
disclosed so that a reader cannot mistake the ceiling for a second reported
architecture.

For numerical binding, eight retained Conv3x3/ConvTranspose weights are
quantized per output channel to narrow-range signed INT8 with a power-of-two
scale.  An independently checked export has zero pre-clip violations; the
largest absolute-code accumulator bound for C1 is 200{,}219, below Acc24.
An operator-local replay covers 160 full calls and $3.27168\times10^9$ output
elements.  Against FP32 it gives 0.009106 MAE, 0.012407 RMSE, 0.239204
maximum absolute error, and 0.9999789 cosine similarity.  All 1{,}280
independently sampled Python-integer probes match; the observed final
accumulator range is $[-29{,}680,27{,}619]$.  These are operator-local
numerical results, not task AEE or a hardware speedup.

\begin{table}[t]
\caption{Frozen-population local DSEC accuracy against a historical,
backend-mismatched baseline (825 frames; lower is better).}
\label{tab:accuracy}
\centering
\scriptsize
\begin{tabular}{@{}lrrr@{}}
\toprule
Metric & Baseline & Candidate & $\Delta$ \\
\midrule
AEE & 1.199514 & 1.197367 & $-0.002147$ \\
AAE & 5.400641 & 5.412808 & $+0.012167$ \\
AAE-Bench. & 5.106363 & 5.121619 & $+0.015256$ \\
DSEC-Fl & 5.313360 & 5.328834 & $+0.015475$ \\
\bottomrule
\end{tabular}
\end{table}

Table~\ref{tab:accuracy} uses the same 825 frames, 18 sequences, and
48{,}152{,}523 valid pixels.  The candidate applies hardware-order attention
to all 12 blocks and dyadic-INT8 QDQ to four C1 Conv3x3 and four decoder
ConvTranspose weights.  Other operators remain at checkpoint precision.  The
historical baseline enabled TF32/cuDNN benchmarking whereas the candidate
disables both; the table is therefore a compatibility gate, not a causal
accuracy improvement.  Aggregate AEE falls by 0.002147, but the three
auxiliary errors rise and 10 of 18 sequences regress in AEE.

\begin{table}[t]
\caption{Admitted C1 evidence; open claims are not inferred.}
\label{tab:main}
\centering
\scriptsize
\begin{tabular}{@{}L{0.14\columnwidth}L{0.36\columnwidth}L{0.36\columnwidth}@{}}
\toprule
Class & Result & Boundary \\
\midrule
Cycle model & 1.6945$\times$; $-40.99\%$ time vs.\ strong zero\modeltag & Ten-sample same-ledger ratio; not an RTL speedup \\
Directed RTL & 64-row tile: 196 issues, 58 parents, 31 dead writes, 64 commits\vcstag & Event-count calibration; synthetic signed-12-bit lanes \\
Physical & 166{,}514~$\mu$m$^2$; $+27.9$/$+1.8$~ps; 16{,}549 FM pts; 29.08~mW / 22.07~nJ\dctag & Prelayout nine-SRAM island; mixed-corner window; no SPEF \\
\bottomrule
\end{tabular}
\end{table}

\section{Related Work}
Prosperity isolates product sparsity from bit sparsity on PTB~\cite{prosperity};
the latter comparison is the relevant mechanism-level analogue to
Table~\ref{tab:c1ladder}.  Phi stores hierarchical patterns chosen by
DSE~\cite{phi}.  We adopt their evaluation pattern---cycle simulation plus
RTL/physical anchors and a baseline ladder---without relabeling their numbers
as ours.  FireFly-T motivates multi-spike service in a different
operator~\cite{fireflyt} and is not an iso-graph baseline for C1.

SpinalFlow sorts compressed spike streams to preserve state locality, and SNE
makes convolutional work proportional to the incoming event
stream~\cite{spinalflow,sne}.  Both establish event-driven execution.  Neither
removes repeated products under a finite 1RW parent lifetime with atomic
completion.  ODIN is a 28-nm digital neuromorphic core with online
learning~\cite{frenkel2019odin}; it is a silicon neuron/synapse fabric, not a
product-capture convolution island.  TCAS-II SRAM compute-in-memory briefs
push sparsity into the array, including dynamic-sparsity control and PIM
tutorials~\cite{qiao2022dscim,taylor2022pim}.  C1 is a digital 1RW parent store
with residual reconstruction, not an in-array MAC macro; we cite those briefs
as the circuits-venue prior on sparsity, not as interchangeable PPA.

A 28-nm optical-flow chip jointly reports operations, external-memory
traffic, latency, energy, and AEE~\cite{ciccflow}.  We follow that separation
of metrics and do not numerically rank its silicon measurements against this
prelayout island.  For a co-designed network, a published chip rarely supports
the identical graph.  Comparison therefore stays at three levels: same-network
component baselines, iso-workload public-artifact opportunities, and published
technology-normalized metrics.  Results from different levels are never multiplied
into a system speedup.

\section{Discussion}
C1's circuits contribution is the 1RW contract, not a larger
sparsity percentage.  On this ledger, bit sparsity is already
visible to a zero-skipping baseline and does not move time
(1.003$\times$).  Product reuse does, but only while a legal parent
occupies one of nine macros and the arbiter can still meet the
row deadline.  Removing either constraint returns the design to
the 1.902$\times$ ceiling, which we refuse to report as an island.

ELSA and SpikeX reuse bundled events or weights across
windows~\cite{elsa,spikex}.  Those priors matter for a typed
source stream, but they do not substitute for residual
reconstruction of a convolution product under a live-parent
bitmap.  Claiming them as C1 would mix delivery reuse with
product reuse.  This brief keeps the two objects separate: only
product reuse with atomic completion is measured here.

The physical numbers are prelayout on purpose.  ZeroWireload
and ideal clocks are optimistic for a nine-macro island; they
are still the correct class of evidence for an Express Brief that
does not have a taped-out chip.  A reader looking for pJ/SOP
on silicon should use ODIN or the CICC optical-flow
chip~\cite{frenkel2019odin,ciccflow} and should not scale our
29.08~mW window to those metrics.

The ten-sample spread (1.661--1.723$\times$) is narrow enough to
argue that C1 is not a one-trace accident, and wide enough to
forbid quoting only the best sample.  We therefore report the
ledger aggregate 1.6945$\times$ and the geometric mean 1.6946$\times$
together, and we do not promote a max-sample headline.  The
directed 64-row VCS tile sits at the other end of that honesty
rule: it is one real-mask slice used to count grants, not a second
speedup.

Deadline holds and stalls in that tile (six and fourteen) are
evidence that the 1RW arbiter is doing work.  A design that
never stalls is either dual-ported or under-occupied.  C1 is
neither.  The 12.25\% tax versus the concurrent-access ceiling
is the same phenomenon at ledger scale: the port, not the
parent relation, is the remaining bottleneck once dead writes
are gone.

\section{Limitations}
The 1.6945$\times$ ratio is a cycle-model result calibrated by one real-mask
RTL tile.  The mapped nine-SRAM component is a separate physical anchor, not
a full integration of the 240-KiB-class common-charge ledger.  Energy is a
directed window, not energy per frame, and the activity factor of that window
is the directed 64-row mask task rather than a DSEC frame.  Hold and power of
a companion weight-broadcast fabric, full-network FPS, and monolithic
integration remain open and are not derived from C1 ratios.

This brief therefore refuses several claims that would make a stronger
headline.  We do not multiply the C1 model ratio by any FC-delivery or
neuron-service factor.  We do not treat the concurrent-access ceiling as a
second architecture.  We do not report Yosys or educational-library area as
ASIC PPA.  We do not promote the compatibility-gate AEE delta as a
quantization gain.  Those omissions are completeness requirements for a
binary Express Brief, not missing experiments to be inferred by the reader.
If a later campaign closes matched hold and power on the weight-broadcast
fabric, that result belongs in a separate manuscript with at least thirty
percent new material relative to this brief; it is not an unstated extra
column of Table~\ref{tab:main}.

\section{Conclusion}
The principal circuit result is not firing sparsity itself, but its
executable capture under a finite single-port parent store.  C1 provides
bounded exact product reuse with a same-ledger model reduction of 1.6945$\times$
and a 28-nm SRAM island that meets 3-ns setup/hold.  Separating model, VCS,
and physical evidence yields a complete Express Brief without an unsupported
whole-network claim.

\bibliographystyle{IEEEtran}
\bibliography{references}
\end{document}
